\documentclass[10pt,a4paper]{amsart}
\usepackage{amsmath,amssymb,amsthm,graphicx}
% Retained preamble for the Stacks Project chapter {CH} (shared project preamble).
% The chapter file pinned-original-source.tex only carries \input{{preamble}} on its first line,
% so no preamble text lives inside the chapter itself.  The Stacks Project uses one single shared
% preamble for every chapter: the snapshot stages it three times (sources/stacks-chow/preamble.tex,
% sources/stacks-etale-cohomology/preamble.tex and focused-seeds/AGP/sources/stacks-snapshot-20261011/
% preamble.tex) and those copies are identical in size.  The block below is copied mechanically,
% line for line, from sources/stacks-chow/preamble.tex lines 14-18 and 154-217: the Xy-pic setup and
% the native theorem declarations together with the mathematical macros.  Nothing was added, removed
% or redefined; the wrapper class and the amsmath/amssymb/amsthm/graphicx base packages are supplied
% by the assembler.
% For commutative diagrams we use Xy-pic
% Chapter: stacks-spaces-more-morphisms.tex -- More on Morphisms of Spaces.
\usepackage[all]{xy}

% We use 2cell for 2-commutative diagrams.
\xyoption{2cell}

% Theorem environments.
%
\theoremstyle{plain}
\newtheorem{theorem}[subsection]{Theorem}
\newtheorem{proposition}[subsection]{Proposition}
\newtheorem{lemma}[subsection]{Lemma}

\theoremstyle{definition}
\newtheorem{definition}[subsection]{Definition}
\newtheorem{example}[subsection]{Example}
\newtheorem{exercise}[subsection]{Exercise}
\newtheorem{situation}[subsection]{Situation}

\theoremstyle{remark}
\newtheorem{remark}[subsection]{Remark}
\newtheorem{remarks}[subsection]{Remarks}

\numberwithin{equation}{subsection}

% Macros
%
\def\lim{\mathop{\mathrm{lim}}\nolimits}
\def\colim{\mathop{\mathrm{colim}}\nolimits}
\def\Spec{\mathop{\mathrm{Spec}}}
\def\Hom{\mathop{\mathrm{Hom}}\nolimits}
\def\Ext{\mathop{\mathrm{Ext}}\nolimits}
\def\SheafHom{\mathop{\mathcal{H}\!\mathit{om}}\nolimits}
\def\SheafExt{\mathop{\mathcal{E}\!\mathit{xt}}\nolimits}
\def\Sch{\mathit{Sch}}
\def\Mor{\mathop{\mathrm{Mor}}\nolimits}
\def\Ob{\mathop{\mathrm{Ob}}\nolimits}
\def\Sh{\mathop{\mathit{Sh}}\nolimits}
\def\NL{\mathop{N\!L}\nolimits}
\def\CH{\mathop{\mathrm{CH}}\nolimits}
\def\proetale{{pro\text{-}\acute{e}tale}}
\def\etale{{\acute{e}tale}}
\def\QCoh{\mathit{QCoh}}
\def\Ker{\mathop{\mathrm{Ker}}}
\def\Im{\mathop{\mathrm{Im}}}
\def\Coker{\mathop{\mathrm{Coker}}}
\def\Coim{\mathop{\mathrm{Coim}}}

% Boxtimes
%
\DeclareMathSymbol{\boxtimes}{\mathbin}{AMSa}{"02}

%
% Macros for moduli stacks/spaces
%
\def\QCohstack{\mathcal{QC}\!\mathit{oh}}
\def\Cohstack{\mathcal{C}\!\mathit{oh}}
\def\Spacesstack{\mathcal{S}\!\mathit{paces}}
\def\Quotfunctor{\mathrm{Quot}}
\def\Hilbfunctor{\mathrm{Hilb}}
\def\Curvesstack{\mathcal{C}\!\mathit{urves}}
\def\Polarizedstack{\mathcal{P}\!\mathit{olarized}}
\def\Complexesstack{\mathcal{C}\!\mathit{omplexes}}
% \Pic is the operator that assigns to X its picard group, usage \Pic(X)
% \Picardstack_{X/B} denotes the Picard stack of X over B
% \Picardfunctor_{X/B} denotes the Picard functor of X over B
\def\Pic{\mathop{\mathrm{Pic}}\nolimits}
\def\Picardstack{\mathcal{P}\!\mathit{ic}}
\def\Picardfunctor{\mathrm{Pic}}
\def\Deformationcategory{\mathcal{D}\!\mathit{ef}}
\newenvironment{reference}{}{}
\newenvironment{slogan}{}{}
\newenvironment{history}{}{}
\usepackage{cleveref}
\crefname{originallemma}{original lemma}{original lemmas}
\Crefname{originallemma}{Original lemma}{Original lemmas}
\makeatletter
\newlabel{flat-lemma-flat-finite-presentation-dimension-over-dense-open}{{081L}{source}{}{sourcepointer-1}{}}
\newlabel{flat-lemma-flat-finite-presentation-dimension-over-dense-open@cref}{{[originallemma][0][]081L}{[1][0][]source}{}{}{}}
\makeatother
\crefname{originallemma}{original lemma}{original lemmas}
\Crefname{originallemma}{Original lemma}{Original lemmas}
\makeatletter
\newlabel{spaces-morphisms-lemma-quasi-compact-immersion}{{088G}{source}{}{sourcepointer-2}{}}
\newlabel{spaces-morphisms-lemma-quasi-compact-immersion@cref}{{[originallemma][0][]088G}{[1][0][]source}{}{}{}}
\makeatother
\crefname{originallemma}{original lemma}{original lemmas}
\Crefname{originallemma}{Original lemma}{Original lemmas}
\makeatletter
\newlabel{spaces-morphisms-section-relative-dimension}{{04NH}{source}{}{sourcepointer-3}{}}
\newlabel{spaces-morphisms-section-relative-dimension@cref}{{[originallemma][0][]04NH}{[1][0][]source}{}{}{}}
\makeatother
\crefname{originallemma}{original lemma}{original lemmas}
\Crefname{originallemma}{Original lemma}{Original lemmas}
\makeatletter
\newlabel{spaces-morphisms-section-scheme-theoretic-closure}{{0831}{source}{}{sourcepointer-4}{}}
\newlabel{spaces-morphisms-section-scheme-theoretic-closure@cref}{{[originallemma][0][]0831}{[1][0][]source}{}{}{}}
\makeatother
\title{Relative dimension bounds for flat maps propagate from a dense open of the base}
\author{The Stacks Project Authors}
\date{stacks-project:spaces-more-morphisms.tex:local-snapshot-20261011; GNU FDL 1.2 or later}
\begin{document}
\maketitle
\noindent\textit{Editorial note.} Statement span is lines 8590-8601 of sources/stacks-spaces-more-morphisms/pinned-original-source.tex; proof span is lines 8603-8622, the complete author proof (18 body lines).\par\medskip
\noindent\textit{Editorial note.} Navigation stubs: 4 cited-but-not-included labels are supplied as non-load-bearing pointers with their pinned tags where the snapshot tag index has them.\par\medskip
\noindent\textit{Editorial note.} extra\_packages is empty on purpose: the retained preamble already loads xy with options and a second declaration clashes in the assembler wrapper.\par\medskip
\noindent\textit{Editorial note.} Text is retained verbatim, including the author's own typos and placeholders.\par\medskip
\noindent\textit{Editorial note.} One of the chapter's dimension-propagation family; the hypothesis split (finite presentation or quasi-compact open) is retained exactly.\par\medskip
\noindent\textit{Allowed standard prerequisites (editorial).}\begin{itemize}
\item Explicit prerequisite boundary. Scheme-theoretic density of quasi-compact dense opens under reduction (Morphisms of Spaces), etale-local character of relative dimension (same chapter's definition section), and the schemes propagation theorem (More on Flatness, lemma-flat-finite-presentation-dimension-over-dense-open).
\end{itemize}\medskip
\ifdefined\hypertarget\hypertarget{sourcepointer-1}{}\fi
\noindent\textit{Source reference.} Original lemma 081L: sources/stacks-spaces-more-morphisms/pinned-original-source.tex, source lines 8622--8622. Cited by the retained statement or proof but not part of the collected spans; supplied only as a navigation pointer. The argument does not depend on its content.\par\medskip
\ifdefined\hypertarget\hypertarget{sourcepointer-2}{}\fi
\noindent\textit{Source reference.} Original lemma 088G: sources/stacks-spaces-more-morphisms/pinned-original-source.tex, source lines 8622--8622. Cited by the retained statement or proof but not part of the collected spans; supplied only as a navigation pointer. The argument does not depend on its content.\par\medskip
\ifdefined\hypertarget\hypertarget{sourcepointer-3}{}\fi
\noindent\textit{Source reference.} Original lemma 04NH: sources/stacks-spaces-more-morphisms/pinned-original-source.tex, source lines 8622--8622. Cited by the retained statement or proof but not part of the collected spans; supplied only as a navigation pointer. The argument does not depend on its content.\par\medskip
\ifdefined\hypertarget\hypertarget{sourcepointer-4}{}\fi
\noindent\textit{Source reference.} Original lemma 0831: sources/stacks-spaces-more-morphisms/pinned-original-source.tex, source lines 8622--8622. Cited by the retained statement or proof but not part of the collected spans; supplied only as a navigation pointer. The argument does not depend on its content.\par\medskip
\begin{lemma}
\label{lemma-flat-fp-dimension-over-dense-open}
Let $S$ be a scheme. Let $f : X \to Y$ be a morphism of algebraic spaces
over $S$ which is flat and locally of finite type. Let $V \subset Y$ be an
open subspace such that $|V| \subset |Y|$ is dense and such that $X_V \to V$
has relative dimension $\leq d$. If also either
\begin{enumerate}
\item $f$ is locally of finite presentation, or
\item $V \to Y$ is quasi-compact,
\end{enumerate}
then $f : X \to Y$ has relative dimension $\leq d$.
\end{lemma}

\begin{proof}
We may replace $Y$ by its reduction, hence we may assume $Y$ is reduced.
Then $V$ is scheme theoretically dense in $Y$, see
Morphisms of Spaces, Lemma
\ref{spaces-morphisms-lemma-quasi-compact-immersion}.
By definition the property of having relative dimension $\leq d$ can
be checked on an \'etale covering, see
Morphisms of Spaces, Sections \ref{spaces-morphisms-section-relative-dimension}.
Thus it suffices to prove $f$ has relative dimension $\leq d$
after replacing $X$ by a scheme surjective and \'etale over $X$.
Similarly, we can replace $Y$ by a scheme surjective and
\'etale and over $Y$. The inverse image of $V$ in this scheme is scheme
theoretically dense, see
Morphisms of Spaces, Section
\ref{spaces-morphisms-section-scheme-theoretic-closure}.
Since a scheme theoretically dense open of a scheme is in particular
dense, we reduce to the case of schemes which is
More on Flatness, Lemma
\ref{flat-lemma-flat-finite-presentation-dimension-over-dense-open}.
\end{proof}

\end{document}
